% Original PDF page 26; hash in recovery-map.json.
\recoverypage{26}
\section*{Q  Additional retrieval and proof-assistance evaluation}
\noindent
The training arms T0--T5 remain unchanged. Retrieval and proof-assistance conditions should be\linebreak[4]
evaluated on the same frozen tasks, with controlled budgets, rather than folded into an unexplained\linebreak[4]
end-to-end score. The following are proposed experiments; no outcomes are claimed in this revision.\par

% Editable layout transcription; cell boundaries not asserted.
\begin{center}\resizebox{\linewidth}{!}{\begin{tabular}{@{}l@{}}
\hspace*{11\fontdimen2\font}Table 13: Separate measurements for retrieval, planning, and proof assistance. \\
 \\
Comparison\hspace*{25\fontdimen2\font}Controlled variable\hspace*{16\fontdimen2\font}Required outcome \\
Lexical / vector / graph\hspace*{11\fontdimen2\font}Same corpus, queries, target\hspace*{7\fontdimen2\font}Independently judged relevant \\
\hspace*{35\fontdimen2\font}interpretation and budget\hspace*{10\fontdimen2\font}context and accepted-premise \\
\hspace*{70\fontdimen2\font}yield \\
Baseline / graph selector\hspace*{10\fontdimen2\font}Same goal and candidate\hspace*{12\fontdimen2\font}Recall against stated labels; \\
\hspace*{35\fontdimen2\font}corpus\hspace*{29\fontdimen2\font}proxy status; actual useful \\
\hspace*{70\fontdimen2\font}proof coverage \\
Unguided / guided plan\hspace*{13\fontdimen2\font}Same admitted sources and\hspace*{10\fontdimen2\font}Deterministic replay, subgoal \\
\hspace*{35\fontdimen2\font}route budget\hspace*{23\fontdimen2\font}coverage, abstention and total \\
\hspace*{70\fontdimen2\font}cost \\
Hammer / Leanstral assistance\hspace*{6\fontdimen2\font}Same source goal, native\hspace*{11\fontdimen2\font}Candidate rate, reconstruction \\
\hspace*{35\fontdimen2\font}checker and resource envelope\hspace*{6\fontdimen2\font}acceptance, unsupported and \\
\hspace*{70\fontdimen2\font}timeout rates \\
Proof heads off / on\hspace*{15\fontdimen2\font}Same representation checkpoint\hspace*{5\fontdimen2\font}Calibration, route value, \\
\hspace*{35\fontdimen2\font}and isolated update policy\hspace*{9\fontdimen2\font}protected-objective invariance \\
Guidance off / promoted\hspace*{12\fontdimen2\font}Same compiler profile and\hspace*{10\fontdimen2\font}Actual consumer activation, \\
\hspace*{35\fontdimen2\font}frozen source families\hspace*{13\fontdimen2\font}independent source fidelity, \\
\hspace*{70\fontdimen2\font}regressions
\end{tabular}}\end{center}

\noindent
For all conditions, retain the corpus and index revision, tokenizer, embedding producer, index\linebreak[4]
implementation, model digest, admitted premises, goal root, logical profile, supported translation\linebreak[4]
fragment, toolchain, checker, time/memory limits, and outcome. Record fixture or mock targets\linebreak[4]
explicitly. Count unsupported, unavailable, and abstained cases in coverage denominators rather than\linebreak[4]
reporting only successful requests. End-to-end cost includes target construction, indexing, gradient\linebreak[4]
updates, model calls, failed candidates, proof reconstruction, semantic validation, and human review.\linebreak[4]
These experiments distinguish three hypotheses: better access to relevant context; better search under\linebreak[4]
a fixed formal interpretation; and better fidelity of that interpretation to the source. Improvements in\linebreak[4]
one do not establish improvements in the other two. No new performance value, trained checkpoint,\linebreak[4]
native solver execution, or deployment qualification is inferred from rebuilding this paper.\par

